% The source was compiled and checked against the supplied 28-page PDF.
\documentclass[11pt]{amsart}

\usepackage[margin=1.15in]{geometry}
\usepackage[T1]{fontenc}
\usepackage{lmodern}
\usepackage{amssymb,mathtools}
\usepackage{microtype}
\usepackage{enumitem}
\usepackage{booktabs,tabularx,placeins}
\usepackage{xcolor}
\usepackage[pagebackref=true,colorlinks=true,linkcolor=blue!55!black,citecolor=blue!55!black,urlcolor=blue!55!black]{hyperref}
\renewcommand*{\backref}[1]{}
\renewcommand*{\backrefalt}[4]{%
  \ifnum#1>0\relax\ensuremath{\uparrow}\,#2\fi}
\renewcommand*{\backrefsep}{,~}
\renewcommand*{\backreftwosep}{,~}
\renewcommand*{\backreflastsep}{,~}

\newtheorem{theorem}{Theorem}[section]
\newtheorem{proposition}[theorem]{Proposition}
\newtheorem{lemma}[theorem]{Lemma}
\newtheorem{corollary}[theorem]{Corollary}
\theoremstyle{definition}
\newtheorem{definition}[theorem]{Definition}
\theoremstyle{remark}
\newtheorem{remark}[theorem]{Remark}
\numberwithin{equation}{section}

\newcommand{\C}{\mathbb C}
\newcommand{\R}{\mathbb R}
\newcommand{\PP}{\mathbf P}
\newcommand{\OO}{\mathcal O}
\newcommand{\Tot}{\operatorname{Tot}}
\newcommand{\Bl}{\operatorname{Bl}}
\newcommand{\ddc}{dd^c}
\newcommand{\psh}{\mathrm{psh}}
\newcommand{\Levi}{\mathcal L}
\newcommand{\barp}{\bar\partial}
\newcommand{\supp}{\operatorname{supp}}
\newcommand{\base}{\operatorname{pr}}
\newcommand{\Uone}{\mathrm U(1)}

\title[Continuous and smooth plurisubharmonic exhaustions]{Continuous and Smooth
Plurisubharmonic Exhaustions:\\
Threefold Counterexamples and Surface Regularization}
\author{Zhijian Chen}
\address{Institute of Mathematics, Academy of Mathematics and Systems Science, Chinese Academy of Sciences, Beijing 100190, China}
\email{chenzhijian@amss.ac.cn}
\author{Bo Xiao}
\address{School of Mathematical Sciences, University of Chinese Academy of Sciences, Beijing, China}
\email{xb1994@amss.ac.cn}
\date{October 29, 2026}
\subjclass[2020]{Primary 32F17; Secondary 32Q15, 32L05, 32U40}
\keywords{Plurisubharmonic exhaustion, K\"ahler manifold, semipositive line bundle,
ruled surface, Hironaka modification, Ueda theory, class VII surface, boundary divisor, positive current}

\setcounter{tocdepth}{1}
\begin{document}

\begin{abstract}
We construct smooth K\"ahler and non-K\"ahler threefolds with locally
Lipschitz plurisubharmonic exhaustions but no $C^2$ plurisubharmonic
exhaustion. The K\"ahler example is quasi-projective and arises from
a continuous semipositive metric on Koike's modified Grauert example;
a radius criterion excludes every $C^2$ exhaustion. In dimension two,
we prove continuous-to-smooth regularization for divisor complements
in compact K\"ahler surfaces, hence for smooth quasi-projective
surfaces. The proof combines level averaging of positive
$dd^c$-closed currents with the surface current theorem and Hodge
index. We also give boundary existence criteria and prove holomorphic
convexity for reduced pure $n$-dimensional spaces with a continuous
psh exhaustion and a holomorphic map to $\C^{n-1}$ whose fibers all
have dimension at most one. The general noncompact surface problem
remains open here.
\end{abstract}
\maketitle
\tableofcontents
\newpage

\section{Introduction and main results}\label{sec:intro}

A continuous plurisubharmonic exhaustion need not admit a smooth
replacement, even on a quasi-projective manifold. We construct
threefold examples and prove affirmative results for compactifiable
surfaces and spaces fibered by curves. Throughout, psh means
plurisubharmonic and smooth means $C^\infty$.

\begin{theorem}\label{thm:main}
Let $C$ be the smooth projective completion of $y^2=x^5-x$, let
$K=K_C$, and choose a nonsplit extension
\begin{equation}\label{eq:extension-intro}
0\longrightarrow K\longrightarrow V\longrightarrow\OO_C
\longrightarrow0.
\end{equation}
Set
\begin{equation}\label{eq:main-line}
\begin{gathered}
S=\PP_{\mathrm{lines}}(V),\quad \pi:S\to C,\quad Y=\PP(K),\\
L=\OO_S(Y)\otimes\pi^*K,\qquad M=\Tot(L^{-1}).
\end{gathered}
\end{equation}
Then $S$ is projective and $M$ is a smooth quasi-projective K\"ahler
threefold. The line bundle $L$ admits a locally Lipschitz Hermitian
metric with psh weights, but no $C^2$ semipositive Hermitian metric.
The threefold $M$ admits a locally Lipschitz psh exhaustion, but no
$C^2$ psh exhaustion.
\end{theorem}

The ruled surface is Koike's modified Grauert example
\cite[Example~4.1]{Koike2021}; the semipositivity obstruction is known
\cite[Theorem~1.4]{Koike2021}\cite[Remark~2.8]{Wu2026}.
Our proof supplies an explicit continuous metric and uses the radius
criterion (Corollary~\ref{cor:criterion}): for a line bundle $A$ over
a compact complex manifold and $r\in\{2,\infty\}$,
\[
\begin{gathered}
\Tot(A^{-1})\text{ has a }C^r\text{ psh exhaustion}
\\\Longleftrightarrow\\
A\text{ has a }C^r\text{ semipositive metric}.
\end{gathered}
\]
The obstruction therefore applies to every $C^2$ exhaustion. The
$C^1$ threshold remains undetermined.

\begin{theorem}\label{thm:nonkahler-threefold}
For every $x\in M$ there is a smooth non-K\"ahler threefold $N$ and a
proper holomorphic modification $\mu:N\to M$ such that
\[
N\setminus\mu^{-1}(x)\simeq M\setminus\{x\}.
\]
The threefold $N$ admits a locally Lipschitz psh exhaustion, but no
$C^2$ psh exhaustion.
\end{theorem}

Hironaka's opposite-order blowups produce the non-K\"ahler
configuration over $x$. Proposition~\ref{prop:compact-modification}
transfers the exhaustion properties through this compactly supported
modification.

\begin{theorem}\label{thm:surface-introduction}
Let $S$ be a smooth compact connected K\"ahler surface and $D\ne0$
an effective divisor. If $S\setminus\supp D$ admits a continuous
psh exhaustion, it admits a smooth one. Consequently, the same
implication holds for every smooth complex quasi-projective surface.
\end{theorem}

The proof averages positive $dd^c$-closed currents over exhaustion
levels to obtain a current with zero Lelong numbers. The current
theorem of Dinh--Nguy\^en--Sibony \cite{DNS22} and Hodge index then
give the following alternative: either $S\setminus\supp D$ is a
modification of a Stein space, or a divisor
$A=\sum a_iD_i$, $a_i\in\mathbb Z_{>0}$, satisfies
\[
\supp A=\supp D,\qquad A^2=0,\qquad
\OO_S(A)\text{ is unitary flat near }D.
\]
Its canonical section yields a pluriharmonic logarithmic collar.
No normal-crossing or minimality hypothesis is imposed on $S,D$.
Section~\ref{sec:surfaces} also gives separate numerical and
neighborhood criteria for existence of an exhaustion.

\begin{theorem}\label{thm:curve-fiber-introduction}
Let $X$ be a reduced complex space of pure dimension $n\ge2$ with a
continuous real-valued psh exhaustion. If a holomorphic map
\[
F:X\longrightarrow\C^{n-1}
\]
has every nonempty fiber of dimension at most one, then $X$ is
holomorphically convex and admits a smooth psh exhaustion.
\end{theorem}

The fiber bound includes exceptional fibers. Neither properness of
$F$ nor normality of $X$ nor a restriction on $\operatorname{Sing}X$
is required. Analytic fiber cycles and Bishop compactness establish
the stability needed for V\^{\i}j\^{\i}itu's convexity criteria
\cite{VijituForum2022}. For irreducible surfaces, convexity already
follows from \cite[Theorem~1.1]{Vajaitu2022}.

These regularization results assert existence of a smooth exhaustion,
not approximation of the prescribed one. They do not settle the
problem for arbitrary noncompact complex surfaces. Within smooth
quasi-projective manifolds, Theorems~\ref{thm:main}
and~\ref{thm:surface-introduction} place the least possible
counterexample dimension at three. The separate elliptic-neighborhood
criterion uses the September 2026 preprint \cite{CFSTZ2026}; the main
surface theorem does not depend on that input.

\paragraph{Conventions.}\label{sec:conventions}
We use
\[
d^c=\frac{i}{4\pi}(\barp-\partial),\qquad
\ddc=\frac{i}{2\pi}\partial\barp,\qquad
\|e\|_h^2=e^{-\varphi},\qquad\Theta_h(A)=\ddc\varphi.
\]
For local frames $e_j=g_{ij}e_i$,
\begin{equation}\label{eq:metric-transition}
\varphi_j=\varphi_i-\log|g_{ij}|^2.
\end{equation}
A continuous Hermitian metric is semipositive when its weights are
psh. For $C^2$ functions, write
$\Levi_f(\xi)=\sum f_{z_j\bar z_k}\xi_j\overline{\xi_k}$.
For the canonical section $s_D$,
\begin{equation}\label{eq:PL}
\ddc\log\|s_D\|_h^2=[D]-\Theta_h(\OO(D)).
\end{equation}
An exhaustion is a continuous real-valued function with compact
sublevel sets. Continuous or smooth weak completeness means existence
of a psh exhaustion of that regularity. On complex spaces, smoothness
means local restriction of smooth ambient functions.
The projectivization $\PP_{\mathrm{lines}}(E)$ parametrizes lines,
with tautological bundle $\OO_{\PP(E)}(-1)$.

\section{The ruled surface and its metrics}\label{sec:metrics}

\subsection{Geometry and the first-order obstruction}\label{sec:ruled}

The map $C\to\PP^1$ defined by $y^2=x^5-x$ has six simple branch
points, including infinity. Hence
\[
 2g(C)-2=-4+6=2,\qquad \deg K=2,\qquad
 \operatorname{Ext}^1(\OO_C,K)=H^1(C,K)\simeq\C.
\]
Choose a nonzero extension class as in \eqref{eq:extension-intro}.
All such classes give isomorphic middle terms after scalar
automorphisms of the end terms. By GAGA, $V$ is algebraic; thus
$S=\PP_{\mathrm{lines}}(V)$ is projective. For $Y=\PP(K)$ and
$L=\OO_S(Y)\otimes\pi^*K$,
\begin{equation}\label{eq:restriction}
 N_{Y/S}\simeq\operatorname{Hom}(K,V/K)\simeq K^{-1},\qquad
 Y^2=-2,\qquad L|_Y\simeq\OO_Y.
\end{equation}
The last trivialization is induced by $K^{-1}\otimes K$.
If $F$ is a fiber and $D\ne Y$ is an irreducible horizontal curve,
\[
 L\cdot Y=0,\quad L\cdot F=1,\quad L^2=2,\quad
 L\cdot D=Y\cdot D+2\deg(\pi|_D)>0.
\]
Thus $L$ is nef and big. This is the modified Grauert example of
\cite[Example~4.1]{Koike2021}.

Choose local frames $k_i$ of $K$ and lifts $v_i$ of $1\in\OO_C$,
with
\[
 k_j=g_{ij}k_i,\qquad v_j=v_i+c_{ij}k_i.
\]
In the coordinates $[k_i+w_iv_i]$ near $Y$, and the frames
$\ell_i=d_i\otimes\pi^*k_i$ with $s_Y=w_id_i$, one obtains
\begin{equation}\label{eq:L-transition}
 w_j=\frac{g_{ij}w_i}{1-c_{ij}w_i},\qquad
 \ell_j=(1-c_{ij}w_i)\ell_i.
\end{equation}

\begin{proposition}\label{prop:no-metric}
The line bundle $L$ has no $C^2$ semipositive Hermitian metric.
\end{proposition}
\begin{proof}
Suppose its local weights $\varphi_i$ are $C^2$ and psh.
Their restrictions to $Y$ agree and are subharmonic, hence constant.
After rescaling, $\varphi_i(z,0)=0$. Positivity of the Levi matrix and
$\varphi_{i,z\bar z}(z,0)=0$ imply
$\varphi_{i,w\bar z}(z,0)=0$. Thus
$a_i(z)=\varphi_{i,w}(z,0)$ is holomorphic. Differentiating the metric
transition law gives
\[
 \varphi_j(z,w_j)=\varphi_i(z,w_i)-\log|1-c_{ij}w_i|^2
 \quad\Longrightarrow\quad g_{ij}a_j=a_i+c_{ij}.
\]
Consequently $v_i-a_ik_i$ glue to a holomorphic lift of $1$, splitting
the chosen extension. This contradiction is the first-order
obstruction of \cite[Theorem~1.4]{Koike2021}; the calculation explicitly
uses only $C^2$ regularity.
\end{proof}

\subsection{Cusp correction and a global metric}\label{sec:cusp}

\begin{lemma}[Cusp correction]\label{lem:cusp}
Let $D$ be a compact smooth curve in a complex surface, with
$\deg N_{D/X}<0$, and let $A|_D$ be unitary flat. Choose smooth
metrics $h_0$ on $A$ and $h_D$ on $\OO_X(D)$ such that $h_0|_D$ is
flat and $\Theta_{h_D}|_{TD}<0$. Put $q=\|s_D\|_{h_D}^2$.
For $0<a\le\tfrac12$, the local weights
\[
 \varphi_0+\Lambda q^a
\]
are psh near $D$, strictly psh off $D$, and locally $C^{0,2a}$,
for a suitable $\Lambda>0$. One may take $\Lambda=1$ when
$a<\tfrac12$; a sufficiently large $\Lambda$ works when
$a=\tfrac12$.
\end{lemma}
\begin{proof}
The required metrics exist because negative degree gives a negative
curvature metric on $N_{D/X}$, and metrics on restricted bundles
extend by extending their logarithmic ratios to fixed smooth metrics.
In a finite coordinate cover of $D$, write
\[
 D=\{w=0\},\quad q=|w|^2e^{\beta(z,w)},\quad
 \beta_{z\bar z}(z,0)>0,\quad \tau=\log q.
\]
For $w\ne0$ and $v=(\xi,w\nu)$,
\begin{align}\label{eq:scaled-cusp}
 q^{-a}\Levi_{q^a}(v)
 &=a\Levi_\tau(v)+a^2|\partial\tau(v)|^2\notag\\
 &=a\beta_{z\bar z}(z,0)|\xi|^2
   +a^2|\nu+\beta_z(z,0)\xi|^2
   +O(|w|)(|\xi|^2+|\nu|^2).
\end{align}
The leading form is positive definite. Compactness and
$q^a\asymp|w|^{2a}$ therefore give, uniformly near $D$,
\[
 \Levi_{q^a}(\xi,w\nu)\ge c|w|^{2a}(|\xi|^2+|\nu|^2).
\]
Flatness of $h_0|_D$ gives $\varphi_{0,z\bar z}(z,0)=0$;
smoothness then yields
\[
 |\Levi_{\varphi_0}(\xi,w\nu)|
 \le C|w|(|\xi|^2+|\nu|^2).
\]
Hence
\begin{equation}\label{eq:cusp-coefficient}
 \Levi_{\varphi_0+\Lambda q^a}(\xi,w\nu)
 \ge (\Lambda c|w|^{2a}-C|w|)(|\xi|^2+|\nu|^2).
\end{equation}
For $a<\tfrac12$, shrink the neighborhood with $\Lambda=1$;
for $a=\tfrac12$, choose $\Lambda>C/c$. This proves strict positivity
off $D$. Since $q^a=|w|^{2a}e^{a\beta}$ is locally $C^{0,2a}$,
the weights extend continuously across $D$; the removable singularity
theorem makes the extension psh. The correction is globally defined,
so the metric transition law is preserved.
\end{proof}

\label{sec:global}
Choose a positive metric $h_K$ on $K$ and extend $h_K^{-1}$ on
$\OO_S(Y)|_Y$ to a smooth metric $h_Y$ on $\OO_S(Y)$. Set
\begin{equation}\label{eq:baseline-metric}
 h_0=h_Y\otimes\pi^*h_K,\qquad
 \theta_0=\Theta_{h_0}(L),\qquad q=\|s_Y\|_{h_Y}^2.
\end{equation}
Then $h_0|_Y$ is flat and $\Theta_{h_Y}|_{TY}=-\Theta_{h_K}<0$.
Choose $a,\Lambda$ as in Lemma~\ref{lem:cusp}, and
$0<\varepsilon<Q=\max_Sq$ such that
$\varphi_0+\Lambda q^a$ is psh on $\{q<\varepsilon\}$. Let
\begin{equation}\label{eq:chi}
 \begin{split}
 B&>\max_{\varepsilon\le t\le Q}(\Lambda t^a-\log t),\\
 \chi(0)&=0,\qquad
 \chi(t)=\max\{\Lambda t^a,\log t+B\}\quad(t>0),
 \end{split}
\end{equation}
and define
\begin{equation}\label{eq:global-metric}
 h=h_0e^{-\chi(q)}.
\end{equation}

\begin{proposition}\label{prop:continuous-metric}
The metric $h$ is continuous and semipositive, with locally
$C^{0,2a}$ weights. In particular it is locally Lipschitz for
$a=\tfrac12$.
\end{proposition}
\begin{proof}
On $S\setminus Y$, Poincar\'e--Lelong gives
\begin{equation}\label{eq:log-curvature}
 \theta_0+\ddc\log q=\pi^*\Theta_{h_K}(K)\ge0.
\end{equation}
On $\{0<q<\varepsilon\}$ the weights are maxima of the psh cusp
and logarithmic branches. The logarithmic branch strictly dominates
on $\{q\ge\varepsilon\}$ and on a collar of its boundary, so
\eqref{eq:log-curvature} proves positivity there. Near $Y$ only the
cusp branch is active. Thus the weights are continuous and psh
everywhere. They are $C^{0,2a}$ near $Y$ and maxima of smooth
functions away from $Y$.
\end{proof}

\begin{remark}\label{rem:minimal}
For $a=\tfrac12$, $T=\theta_0+\ddc\chi(q)$ has locally Lipschitz
potentials and minimal singularities. Indeed, any positive current
$\theta_0+\ddc v\in c_1(L)$ has a global quasi-psh potential $v$,
and compactness gives $v\le\chi(q)+C_v$. This answers the
bounded-potential question of \cite[Section~4]{Zhang2025} for
Koike's modified example; no identification with Grauert's original
line bundle is asserted.
\end{remark}

\section{Exhaustions and the radius criterion}\label{sec:radius}

\begin{lemma}\label{lem:norm}
If $A\to Z$ is a holomorphic line bundle over a compact complex
manifold and $h$ is a continuous semipositive Hermitian metric, then
$\rho(v)=\|v\|_{h^{-1}}^2$ is a continuous psh exhaustion of
$\Tot(A^{-1})$. It inherits the local H\"older or Lipschitz
regularity of $h$.
\end{lemma}
\begin{proof}
For $v=\lambda_i e_i^*$ and $\|e_i\|_h^2=e^{-\varphi_i}$,
\begin{equation}\label{eq:local-norm}
 \rho=|\lambda_i|^2e^{\varphi_i}.
\end{equation}
Off the zero section, $\log\rho$ is psh, hence so is its exponential.
Continuity and removability extend this across the zero section.
Compactness of $Z$ makes $\rho$ comparable to a fixed smooth squared
norm; its closed sublevels are therefore compact disk subsets.
The regularity follows from \eqref{eq:local-norm}.
\end{proof}

\begin{lemma}[Compact-group averaging]\label{lem:average}
If a compact group $G$ acts continuously on a locally compact
Hausdorff space and $F$ is a continuous exhaustion, then
$\overline F(x)=\int_GF(gx)\,dg$ is an exhaustion.
Averaging over a smooth holomorphic compact Lie group action
preserves $C^r$ plurisubharmonicity, $r\in\{2,\infty\}$.
\end{lemma}
\begin{proof}
For $K_c=\{F\le c\}$,
\[
 \{\overline F\le c\}\subset G\cdot K_c.
\]
Indeed, outside the compact saturation the orbitwise minimum of
$F$ exceeds $c$. The left side is closed, hence compact.
Differentiation under the integral and averaging Levi forms prove
the remaining assertions.
\end{proof}

\begin{proposition}[Radius criterion]\label{prop:radius}
If $\Tot(A^{-1})$ admits a $C^r$ psh exhaustion, where $A$ is a
holomorphic line bundle over a compact complex manifold and
$r\in\{2,\infty\}$, then $A$ admits a $C^r$ semipositive
Hermitian metric.
\end{proposition}
\begin{proof}
Average the exhaustion over fiberwise circle rotations. In a local
frame the resulting invariant function has the form
\[
 F(z,\lambda_i)=U_i(z,t),\quad t=\log|\lambda_i|^2,
 \qquad V_i(z,\zeta)=U_i(z,\operatorname{Re}\zeta)=F(z,e^{\zeta/2}).
\]
The function $V_i$ is psh; thus $U_{i,tt}\ge0$. For each $z$,
\[
 U_i(z,-\infty)=F(z,0),\qquad U_i(z,+\infty)=+\infty.
\]
Convexity implies $U_{i,t}\ge0$, and $U_{i,t}>0$ whenever
$U_i>F(z,0)$: a zero derivative forces constancy on the preceding
half-line. Fix $c>\max_ZF|_Z$. There is a unique $C^r$ function
$\sigma_i$ such that
\begin{equation}\label{eq:implicit-radius}
 U_i(z,\sigma_i(z))=c,\qquad U_{i,t}(z,\sigma_i(z))>0.
\end{equation}
Twice differentiating gives
\begin{align*}
 0={}&U_{i,z_j\bar z_k}+U_{i,z_jt}\sigma_{i,\bar z_k}
 +U_{i,t\bar z_k}\sigma_{i,z_j}
 +U_{i,tt}\sigma_{i,z_j}\sigma_{i,\bar z_k}
 +U_{i,t}\sigma_{i,z_j\bar z_k}.
\end{align*}
Since the $\zeta$ derivatives of $V_i$ introduce factors $1/2$,
\begin{equation}\label{eq:radius-Levi}
 0\le\Levi_{V_i}(\xi,2\partial\sigma_i(\xi))
 =-U_{i,t}\Levi_{\sigma_i}(\xi).
\end{equation}
Thus $\psi_i=-\sigma_i$ is psh. If $e_j=g_{ij}e_i$, then
$\lambda_j=g_{ij}\lambda_i$, whence
\[
 \sigma_j=\sigma_i+\log|g_{ij}|^2,\qquad
 \psi_j=\psi_i-\log|g_{ij}|^2.
\]
The weights $\psi_i$ define the required metric.
\end{proof}

\begin{corollary}\label{cor:criterion}
For $r\in\{2,\infty\}$ and a line bundle $A$ over a compact complex
manifold,
\[
\begin{gathered}
 \Tot(A^{-1})\text{ admits a }C^r\text{ psh exhaustion}
 \\\Longleftrightarrow\\
 A\text{ admits a }C^r\text{ semipositive metric}.
\end{gathered}
\]
\end{corollary}
\begin{proof}
Combine Proposition~\ref{prop:radius} with
\eqref{eq:local-norm}.
\end{proof}

\begin{proof}[Proof of Theorem~\ref{thm:main}]\label{sec:conclusion}
The total space has the projective compactification
\begin{equation}\label{eq:compactification}
 M=\PP_{\mathrm{lines}}(\OO_S\oplus L^{-1})
 \setminus\PP_{\mathrm{lines}}(L^{-1}),\qquad
 v\longmapsto[(1,v)].
\end{equation}
It is therefore smooth, quasi-projective and K\"ahler. With
$p:M\to S$, the metric of Proposition~\ref{prop:continuous-metric}
gives the exhaustion
\begin{equation}\label{eq:explicit-exhaustion}
 \rho(v)=\|v\|_{h_0^{-1}}^2\exp\bigl(\chi(q(p(v)))\bigr).
\end{equation}
Take $a=\tfrac12$ to obtain local Lipschitz regularity.
A $C^2$ psh exhaustion would give a $C^2$ semipositive metric on $L$
by Proposition~\ref{prop:radius}, contradicting
Proposition~\ref{prop:no-metric}.
\end{proof}

\section{A non-K\"ahler modification over one point}\label{sec:nonkahler}

\begin{proposition}\label{prop:compact-modification}
Let $f:X'\to X$ be a proper surjective holomorphic map of nonempty
complex manifolds, biholomorphic outside $f^{-1}(K)$ for some compact
$K\subset X$. For $r\in\{0,2,\infty\}$, $X'$ admits a $C^r$ psh
exhaustion if and only if $X$ does.
\end{proposition}
\begin{proof}
Pullback preserves psh, regularity and compact sublevels.
Conversely, let $F$ exhaust $X'$, choose $K\subset U\Subset X$ and
$c>\max_{f^{-1}(\overline U)}F$, and take a smooth convex
nondecreasing $\chi$ vanishing on $(-\infty,c]$ and tending to
infinity. Then
\[
 G=\chi\circ F\circ f^{-1}\quad\text{on }X\setminus K,
 \qquad G=0\quad\text{on }K
\]
is $C^r$ and psh, since it vanishes on $U$. For each $a\ge0$, choose
$b$ with $\chi(t)\le a\Rightarrow t\le b$. The closed set
\[
 \{G\le a\}\subset K\cup f(\{F\le b\})
\]
is compact. The case $K=\varnothing$ is immediate.
\end{proof}

\begin{proposition}\label{prop:point-modification}
For every point $x$ of a connected smooth complex threefold $X$,
there is a connected smooth non-K\"ahler threefold $N$ and a proper
modification $\mu:N\to X$ that is biholomorphic over
$X\setminus\{x\}$.
\end{proposition}
\begin{proof}
Blow up $x$:
\[
 b:W=\Bl_xX\to X,\qquad E=b^{-1}(x)\simeq\PP^2.
\]
In $E$, the curves
\[
 A=\{w=0\},\qquad B=\{uv=w^2\}
\]
meet transversely at $p=[1:0:0]$ and $q=[0:1:0]$.
Over $W\setminus\{q\}$ blow up $A$ and then the strict transform
of $B$; over $W\setminus\{p\}$ use the reverse order. On the
overlap the centers are disjoint, so the two blowups commute
canonically. Gluing gives $\nu:N\to W$ and $\mu=b\circ\nu$.
All centers are smooth. Points with distinct images separate in
$W$; points with the same image lie in a common Hausdorff chart
over one of the two base opens. Thus $N$ is a smooth Hausdorff
manifold. Properness is local on the target, the charts are connected
and overlap, and all centers lie over $x$.

We verify the non-K\"ahler obstruction, as in
\cite[Section~6]{Barlet2015} and \cite[Tag~08J1]{StacksHironaka}.
Let $D_A,D_B$ be the exceptional surfaces dominating $A,B$, and let
$\alpha,\beta\in H_2(N,\R)$ be their general fiber classes.
These surfaces are compact and smooth; locally they are ruled
surfaces or their blowups at a point. Their fiber divisors are
linearly equivalent because $A\simeq B\simeq\PP^1$.
Over $p$ and $q$ the respective orders give
\[
 \alpha=[R_p]+[T_p],\quad\beta=[T_p],\qquad
 \beta=[R_q]+[T_q],\quad\alpha=[T_q],
\]
where every displayed component is a smooth rational curve of
multiplicity one. To check the multiplicities at $p$, use
$A=\{y=z=0\}$, $B=\{x=z=0\}$. In the first blowup chart $z=yt$,
the first exceptional divisor is $y=0$ and the second center is
$x=t=0$. The second blowup restricts to the surface blowup of
$(x,t)=(0,0)$; the pullback of the fiber $x=0$ consists of its
strict transform and the exceptional curve, each once.
Consequently
\begin{equation}\label{eq:hironaka-zero-cycle}
 [R_p]+[R_q]=0\quad\text{in }H_2(N,\R).
\end{equation}
The relations take place on compact surfaces, so they hold in
ordinary homology even when $N$ is noncompact. A K\"ahler form
would give
\[
 0=\int_{R_p}\omega+\int_{R_q}\omega>0,
\]
a contradiction.
\end{proof}

\begin{proof}[Proof of Theorem~\ref{thm:nonkahler-threefold}]
Apply Proposition~\ref{prop:point-modification} to $M$ at $x$.
The pullback $\rho\circ\mu$ is a locally Lipschitz psh exhaustion:
$\mu$ is locally Lipschitz and proper. A $C^2$ psh exhaustion on
$N$ would descend after truncation by
Proposition~\ref{prop:compact-modification}, contradicting
Theorem~\ref{thm:main}.
\end{proof}

\section{Surface regularization and boundary criteria}\label{sec:surfaces}
Throughout this section, $X=S\setminus\supp D$, where $S$ is a smooth
compact connected complex surface and $D\ne0$ is effective. A
\emph{divergent collar} is a psh function on a punctured neighborhood
of $\supp D$ tending uniformly to $+\infty$ at the boundary.
An exhaustion of $X$ always has this divergence property.

\subsection{Boundary tools}\label{sec:surface-pair-reduction}
\begin{lemma}\label{lem:surface-boundary-tools}
\begin{enumerate}[label=\textup{(\roman*)},leftmargin=2.5em]
\item\label{item:negative-boundary}
If a connected component of $\supp D$ has negative definite
intersection matrix, $X$ has no continuous psh exhaustion.
\item\label{item:collar-gluing}
Smooth divergent psh collars at all connected boundary components
give a smooth psh exhaustion of $X$.
\item\label{item:matrix-alternatives}
For a connected reduced SNC divisor $E=\bigcup E_i$, let $\lambda$
be the largest eigenvalue of $I=(E_i\cdot E_j)$. If $\lambda>0$,
$E$ has a smooth strictly psh divergent collar. If $\lambda=0$,
$I\le0$, $\dim\ker I=1$, and
\begin{equation}\label{eq:surface-null-boundary}
 A=\sum_i a_iE_i,\qquad a_i\in\mathbb Z_{>0},\qquad
 A\cdot E_i=A^2=0.
\end{equation}
When $S$ is projective, $A$ is nef.
\end{enumerate}
\end{lemma}
\begin{proof}
For \textup{(i)}, Grauert's theorem
\cite[\S4, subsection~8(e)]{Grauert1962} contracts the component to
a point $p$ of a normal surface. A continuous psh function on the
punctured neighborhood is bounded above near $p$: choose a local
finite projection to a bidisc with $p$ the only point over $0$.
Its fiberwise maximum is continuous and psh on the punctured bidisc;
across the branch locus this follows by local boundedness and
removability. The maximum on
$\{|z_1|\le r,\ |z_2|=r\}$ bounds the function on vertical discs
with $z_1\ne0$, and continuity gives the same bound at $z_1=0$
away from the origin. This contradicts divergence at $p$.

For \textup{(ii)}, choose disjoint neighborhoods $W_E$ of the
boundary components, with closures inside the collar domains.
Choose $T_E>\sup_{\partial W_E}\Phi_E$ and a smooth convex
nondecreasing $\chi_E$, zero on $(-\infty,T_E]$ and unbounded
above. Extend $\chi_E(\Phi_E)$ by zero outside $W_E$ and sum.
The sum is smooth, psh, and proper.

For \textup{(iii)}, Perron--Frobenius applied to $I+c\mathrm{Id}$
shows that $\lambda$ is simple and has a strictly positive
eigenvector. When $\lambda>0$, rational approximation gives
$A=\sum a_iE_i$ with $a_i>0$ and $A\cdot E_i>0$ for all $i$.
Thus $\OO_S(A)|_E$ is ample. A projective embedding by a power
gives a Fubini--Study metric on this restriction. Extend its weights
locally to $S$ and patch them relative to a reference metric.
The extensions have positive tangential curvature along $E$;
near a node use the embedding sections themselves, whose
differentials span the ambient cotangent space, so the curvature
is positive in both directions there. Patching can be chosen to
keep these metrics near the nodes. If $q=\|s_E\|^2$ for any smooth
metric on $\OO_S(E)$, then $\ddc q$ is positive in the normal
direction at every smooth point of $E$. Adding $Cq$ to the
weights, with $C\gg1$, gives a metric $h$ of positive curvature
near all of $E$. Hence $-\log\|s_A\|_h^2$ is the required collar.
For $\lambda=0$, the positive kernel has a rational generator;
clearing denominators proves \eqref{eq:surface-null-boundary}.
Effectiveness gives $A\cdot C\ge0$ for every curve $C\not\subset E$,
which proves nefness in the projective case.
\end{proof}

Boundary resolution preserves $X$ when the new boundary is the
full reduced inverse image. Interior blowups preserve both exhaustion
properties by Proposition~\ref{prop:compact-modification}. Contracting
a curve meeting but not contained in the boundary need not give a
proper map of complements; no such reduction will be used.

\begin{proposition}\label{prop:surface-nonconstant-function}
A connected noncompact complex surface with a continuous psh
exhaustion and $\OO(X)\ne\C$ has a smooth psh exhaustion.
\end{proposition}
\begin{proof}
V\^{\i}j\^{\i}itu's theorem \cite[Theorem~1.1]{Vajaitu2022} gives
holomorphic convexity. Pull back a smooth strictly psh exhaustion
of the Stein Remmert reduction under its proper reduction map.
\end{proof}

\begin{proposition}\label{cor:curve-total-space}
For a line bundle $A$ on a compact Riemann surface $B$,
\[
 \begin{gathered}
 \Tot(A^{-1})\text{ continuously weakly complete}\\
 \Longleftrightarrow\quad\deg A\ge0\quad\Longleftrightarrow\\
 \Tot(A^{-1})\text{ smoothly weakly complete}.
 \end{gathered}
\]
\end{proposition}
\begin{proof}
For $d=\deg A\ge0$, choose an area form $\omega$ of integral one
and solve $\ddc f=d\omega-\theta_0$, where $\theta_0$ is the
curvature of a reference metric. The resulting metric has curvature
$d\omega$, and Lemma~\ref{lem:norm} applies. If $d<0$, the infinity
section of $\PP_{\mathrm{lines}}(\OO_B\oplus A^{-1})$ has normal
bundle $A$ and negative square, contradicting
Lemma~\ref{lem:surface-boundary-tools}\textup{(i)}.
\end{proof}

\subsection{Current averaging and regularization}
\label{subsec:divisor-regularization}
\begin{theorem}\label{thm:divisor-regularization}
Let $S$ be compact K\"ahler and $D\ne0$ a reduced divisor.
If $X=S\setminus D$ has a continuous psh exhaustion, it has a
smooth one. If $X$ is not a modification of a Stein space, there is
\[
 A=\sum_i a_iD_i,\qquad a_i\in\mathbb Z_{>0},\qquad
 A\cdot D_i=A^2=0,
\]
such that $\OO_S(A)$ has a smooth Hermitian metric with zero
curvature near $D$. Thus the entire boundary has a smooth
pluriharmonic divergent collar.
\end{theorem}

\begin{lemma}\label{lem:level-averaging}
Let $(S,\omega)$ be a compact K\"ahler surface and $X\subset S$
a connected noncompact open set with a continuous psh exhaustion
$\rho$. If $X$ is not a modification of a Stein space, there is
a positive $(1,1)$-current $T$ on $S$ such that
\begin{equation}\label{eq:averaged-current-properties}
 \ddc T=0,\qquad \int_S T\wedge\omega=1,\qquad
 \supp T\Subset X,\qquad \nu(T,x)=0\quad(x\in S).
\end{equation}
\end{lemma}
\begin{proof}
Let $\Sigma_X^0$ be the intersection of the non-strongly-psh loci
of all continuous psh exhaustions. By
\cite[Lemma~2.1]{MongodiSlodkowski2020},
\begin{equation}\label{eq:kernel-halfline}
 (t_0,\infty)\subset\rho(\Sigma_X^0)
\end{equation}
for some $t_0$. For $t>t_0$, no smooth strictly psh function exists
near $K_t=\{\rho=t\}$. Indeed, if $v$ were such a function,
choose $\delta>0$ with $\{t-\delta\le\rho\le t+\delta\}$ contained
in its domain, and normalize $0<v<1$ there. For
$0<\eta<1$, $0<\varepsilon<\eta\delta/2$, set
$g(s)=s+\eta|s-t|$. The function
\[
 \psi=
 \begin{cases}
 \max\{g(\rho),\rho+\varepsilon v\},& |\rho-t|<\delta,\\
 g(\rho),& |\rho-t|\ge\delta
 \end{cases}
\]
is a continuous psh exhaustion, strongly psh near $K_t$, contrary
to \eqref{eq:kernel-halfline}.

In the Banach space of continuous real $(1,1)$-forms on $K_t$,
the linear subspace $\{\ddc f|_{K_t}:f\in C^\infty(S,\R)\}$
is therefore disjoint from the open cone of positive definite forms.
A local function extends by a cutoff, so this global formulation
loses nothing. Hahn--Banach separation gives a nonzero positive
functional annihilating the subspace. As a current, after mass
normalization it satisfies
\begin{equation}\label{eq:level-current}
 T_t\ge0,\qquad \ddc T_t=0,\qquad
 \supp T_t\subset K_t,\qquad \int_S T_t\wedge\omega=1.
\end{equation}
No closedness of the separated subspace is needed, since the
positive cone is open.

Fix $t_0<a<b$ and choose
\[
 t_{n,j}=a+\frac{j-1/2}{n}(b-a),\qquad
 T_n=\frac1n\sum_{j=1}^nT_{t_{n,j}}.
\]
These currents have mass one and support in the fixed compact shell
$K=\rho^{-1}([a,b])\Subset X$. A subsequence converges weakly to a
positive $dd^c$-closed current $T$ of mass one supported in $K$.

Fix a coordinate ball $B(x,2R)\Subset X$, and put
$\beta=\ddc|z|^2$. Lelong--Jensen monotonicity for positive
$dd^c$-closed currents \cite[Appendix~B, (B.1)--(B.2)]{DNS22} gives,
uniformly in $t$, for $0<r<R$,
\begin{equation}\label{eq:uniform-lelong-bound}
 r^{-2}\int_{B(x,r)}T_t\wedge\beta
 \le R^{-2}\int_{B(x,R)}T_t\wedge\beta\le C_R.
\end{equation}
Only the levels in
$[\inf_{B(x,r)}\rho,\sup_{B(x,r)}\rho]$ can contribute. Counting
grid points yields
\begin{equation}\label{eq:grid-oscillation-bound}
 r^{-2}\int_{B(x,r)}T_n\wedge\beta
 \le C_R\left(\frac{\operatorname{osc}_{B(x,r)}\rho}{b-a}
                  +\frac2n\right).
\end{equation}
Multiply $\beta$ by a cutoff equal to one on $B(x,R)$. The
corresponding positive measures converge weakly, and the open-set
inequality gives
\begin{equation}\label{eq:limit-oscillation-bound}
 r^{-2}\int_{B(x,r)}T\wedge\beta
 \le\frac{C_R}{b-a}\operatorname{osc}_{B(x,r)}\rho
 \longrightarrow0\qquad(r\downarrow0).
\end{equation}
Thus $\nu(T,x)=0$ on $X$; outside $X$, $T$ vanishes locally.
The fixed weak limit works for every ball, so no point-dependent
subsequence or measurable selection is required.
\end{proof}

\begin{lemma}\label{lem:surface-current-inequality}
For a positive $dd^c$-closed $(1,1)$-current $T$ on a compact
K\"ahler surface with all Lelong numbers zero, there is a class
$\alpha_T\in H^{1,1}(S,\R)$ satisfying
\begin{equation}\label{eq:current-class-pairing}
 \alpha_T\cdot[\eta]=\int_S T\wedge\eta
 \qquad(d\eta=0),
\end{equation}
for every smooth real $(1,1)$-form $\eta$. Moreover,
\begin{equation}\label{eq:DNS-inequality}
 \alpha_T^2\ge0,\qquad dT\ne0\ \Longrightarrow\ \alpha_T^2>0.
\end{equation}
If $dT=0$, then $\alpha_T=[T]$.
\end{lemma}
\begin{proof}
The decomposition \cite[\S2.1 and Lemma~B.2]{DNS22}
\[
 T=\Omega+\partial R+\bar\partial\overline R
                  +i\partial\bar\partial u
\]
has $\Omega$ smooth, real, and closed. Integration by parts gives
\eqref{eq:current-class-pairing} with $\alpha_T=[\Omega]$;
nondegeneracy of the intersection pairing makes this class canonical.
The density formula \cite[Theorem~2.2(2), Corollary~2.4]{DNS22} gives
\[
 \alpha_T^2=\|\vartheta\|+2E(T),\qquad
 \|\vartheta\|\ge0,\quad E(T)\ge0,
 \quad E(T)=0\Longleftrightarrow dT=0.
\]
Here $\vartheta$ is the positive measure on the diagonal determined
by a tangent current of $T\otimes T$, and
$E(T)=\int_S\bar\partial R\wedge\partial\overline R$ is the
current energy, independent of the chosen decomposition.
Its mass hypotheses hold because the positive-Lelong-number locus
is empty; no foliation hypothesis is used. This proves
\eqref{eq:DNS-inequality}.
\end{proof}

\begin{proof}[Proof of Theorem~\ref{thm:divisor-regularization}]
The Stein-modification case follows by proper pullback of a smooth
psh exhaustion. Otherwise the preceding lemmas give $T$, supported
in $K\Subset X$, and $\alpha=\alpha_T$ with
\begin{equation}\label{eq:alpha-basic}
 \alpha\cdot[\omega]=1,\qquad \alpha^2\ge0,
 \qquad dT\ne0\Longrightarrow\alpha^2>0.
\end{equation}
For each component $D_i$, choose a metric of curvature $\theta_i$
and canonical section $s_i$. Cut off $\log\|s_i\|^2$ to a smooth
function $f_i$ on $S$ agreeing with it near $K$. Poincar\'e--Lelong
and $dd^c$-closedness give
\begin{equation}\label{eq:alpha-boundary-orthogonality}
 \alpha\cdot[D_i]=\int_S T\wedge\theta_i
          =-\langle T,\ddc f_i\rangle=0.
\end{equation}

For a connected boundary component $E$, let $\lambda_E$ be the
largest eigenvalue of its intersection matrix $I_E$.
Lemma~\ref{lem:surface-boundary-tools}\textup{(i)} excludes
$\lambda_E<0$. If $\lambda_E>0$, rational approximation gives an
integral divisor $B$ supported on $E$ with $B^2>0$.
Then $\alpha\cdot[B]=0$, and Hodge index forces $\alpha^2<0$,
contrary to \eqref{eq:alpha-basic}. Therefore $\lambda_E=0$.
Perron--Frobenius gives a one-dimensional positive rational kernel
of $I_E$. Choose a positive integral generator for each $E$ and sum:
\begin{equation}\label{eq:full-support-null-divisor}
 A=\sum_i a_iD_i,\qquad a_i>0,\qquad
 \supp A=D,\qquad A\cdot D_i=A^2=0.
\end{equation}
These matrix arguments do not require smoothness or normal crossings
of the components.

Since $[A]\ne0$ is isotropic, Hodge index makes $[A]^\perp$
negative semidefinite with radical $\R[A]$. Using
\eqref{eq:alpha-basic}--\eqref{eq:alpha-boundary-orthogonality},
we obtain
\begin{equation}\label{eq:closed-current-null-class}
 dT=0,\qquad [T]=\alpha=c[A],\qquad
 c=\left(\int_A\omega\right)^{-1}>0.
\end{equation}
The $dd^c$-lemma for currents now supplies a real distribution $u$ with
\begin{equation}\label{eq:global-log-potential}
 \ddc u=T-c[A].
\end{equation}
On a neighborhood $U$ of $D$ disjoint from $K$, a local defining
function $f_A$ gives $\ddc(u+c\log|f_A|^2)=0$. Elliptic regularity
for pluriharmonic distributions yields
\begin{equation}\label{eq:local-log-collar}
 u=-c\log|f_A|^2+h,
 \qquad h\text{ smooth and pluriharmonic}.
\end{equation}
Thus $u$ is smooth pluriharmonic on $U\setminus D$ and diverges
uniformly at $D$, including singular points, because every $a_i>0$.
Lemma~\ref{lem:surface-boundary-tools}\textup{(ii)} gives the exhaustion.

For the metric statement, write $s_A=f_i e_i$ and set
$b_i=u/c+\log|f_i|^2$ on $U$. These functions extend smoothly
and pluriharmonically across $D$. If $e_j=g_{ij}e_i$, then
$b_j=b_i-\log|g_{ij}|^2$; hence $\|e_i\|^2=e^{-b_i}$ defines
a flat Hermitian metric and $-\log\|s_A\|^2=u/c$.
Patch it with any global smooth metric using a cutoff equal to one
on a smaller neighborhood of $D$. This gives the claimed global
metric with curvature zero near $D$.
\end{proof}

\begin{corollary}\label{cor:quasiprojective-regularization}
Every smooth complex quasi-projective surface with a continuous
psh exhaustion has a smooth psh exhaustion.
\end{corollary}
\begin{proof}
Resolve a projective closure and principalize the complement ideal,
with centers outside $X$ \cite[Theorems~1.0.1--1.0.2]{Wlodarczyk05}.
This realizes $X$ as the complement of an SNC divisor in a smooth
projective surface; the full inverse image includes any isolated
boundary points. Apply Theorem~\ref{thm:divisor-regularization}.
Compact components need only constant functions, and there are
finitely many components.
\end{proof}

The case $\OO(X)=\C$ occurs: take $X=\Tot(A^{-1})$ for non-torsion
$A\in\operatorname{Pic}^0(B)$. Its squared unitary-flat norm is a
smooth psh exhaustion. The degree-$n$ fiberwise Taylor coefficient
of $f\in\OO(X)$ lies in $H^0(B,A^n)=0$ for $n>0$; the constant
coefficient is constant on $B$. Thus $\OO(X)=\C$.

\begin{corollary}\label{cor:square-zero-continuous-collar}
Let $C\subset S$ be a smooth connected square-zero curve in a compact
K\"ahler surface. A continuous divergent psh collar at $C$ implies
a smooth psh exhaustion of $S\setminus C$. If the complement is not
a Stein modification, $\OO_S(C)$ is unitary flat near $C$ and admits
a pluriharmonic divergent collar.
\end{corollary}
\begin{proof}
Take the maximum of the given collar and a sufficiently large
constant near the boundary of a smaller neighborhood, and extend
by that constant. This gives a continuous exhaustion. In
Theorem~\ref{thm:divisor-regularization}, the single-component kernel
generator can be chosen to be $1$, so $A=C$.
\end{proof}

The theorem asserts existence of a new exhaustion, not approximation
of the given one. For reducible $D$, flatness concerns $\OO_S(A)$,
not necessarily $\OO_S(D)$. No divisor compactification is asserted
for an arbitrary noncompact surface.

\subsection{Numerical existence criteria}\label{sec:surface-nonprojective}
\begin{theorem}\label{thm:nonprojective}
For a nonprojective $S$ and an effective $D\ne0$, the following
are equivalent:
\begin{enumerate}[label=\textup{(\roman*)},leftmargin=2.5em]
\item $S\setminus\supp D$ is continuously weakly complete;
\item $S\setminus\supp D$ is smoothly weakly complete;
\item no connected boundary component has negative definite
intersection matrix.
\end{enumerate}
The Hopf case uses \cite[Theorem~1.6]{CFSTZ2026}.
\end{theorem}
\begin{proof}
Only \textup{(iii)}$\Rightarrow$\textup{(ii)} requires proof.
The classification facts used below are recalled in
\cite[\S2]{FN2005}. If $a(S)=1$, the algebraic reduction is a
holomorphic elliptic fibration $f:S\to B$ with connected fibers,
and every curve is vertical. For a fiber $F=\sum m_iC_i$,
\begin{equation}\label{eq:surface-fiber-intersection-identity}
 x^{\mathsf T}Ix=-\sum_{i<j}(C_i\cdot C_j)m_im_j
       \left(\frac{x_i}{m_i}-\frac{x_j}{m_j}\right)^2,
 \qquad I=(C_i\cdot C_j).
\end{equation}
Every proper principal submatrix is negative definite. Hence the
boundary contains each fiber support it meets, and
$X=f^{-1}(B\setminus P)$ for a finite nonempty $P$.
Pull back a smooth subharmonic exhaustion of $B\setminus P$.

Suppose $a(S)=0$. By \cite[Lemma~2.1]{FN2005}, $c_1(L)^2\le0$
for every line bundle $L$. Perron--Frobenius and \textup{(iii)}
therefore give an effective $A$ of full boundary support with
$A^2=0$. At a blowdown $q:T\to T'$, write
$A=q^*A'+bE$. Then
\[
 0=A^2=(A')^2-b^2,\qquad (A')^2\le0,
\]
so $b=0$. On the minimal model $p:S\to S_0$ we obtain
$A=p^*A_0\ne0$, $A_0^2=0$, and
$\supp A=p^{-1}(\supp A_0)$. The map of complements is proper,
so an exhaustion downstairs suffices.

If $S_0$ is K\"ahler, it is a torus or K3 of algebraic dimension
zero. Such a torus has no curves: an irreducible curve would have
square zero and elliptic normalization, hence be a translate of an
elliptic subtorus and give a nonconstant meromorphic function by the
quotient. On a K3, Riemann--Roch gives
$h^0(\OO(A_0))\ge2$. Both contradict $a(S_0)=0$.

Otherwise $S_0$ is of class~VII. If $b_2(S_0)=n>0$, Enoki's
square-zero-divisor theorem \cite[Main Theorem, p.~454]{Enoki}
makes $A_0$ a multiple of the distinguished cycle of an Enoki surface.
Its complement is an affine-line bundle of degree $-n$ over an
elliptic curve \cite[(0.4), p.~453]{Enoki}. The infinity section of
its $\PP^1$-bundle completion has square $n>0$, so
Lemma~\ref{lem:surface-boundary-tools} gives an exhaustion.
If $b_2(S_0)=0$, a minimal class~VII surface containing a curve is
elliptic or Hopf \cite[p.~453]{Enoki}; $a(S_0)=0$ leaves Hopf.
Its curves are disjoint smooth elliptic curves: adjunction gives
arithmetic genus one, and rational normalizations are impossible
because maps $\PP^1\to S_0$ lift to $\C^2\setminus\{0\}$.
Finite Ueda type gives a smooth strictly psh collar
\cite[Theorem~2.5]{CLPT2018}. Infinite type gives a unitary flat
neighborhood by \cite[Theorem~1.6]{CFSTZ2026}, hence a
pluriharmonic collar. Glue the collars.
\end{proof}

\begin{theorem}\label{cor:ruled-all-degrees}
Let $f:S\to B$ be a smooth $\PP^1$-bundle over a compact Riemann
surface and $D\ne0$ reduced SNC. The three conditions in
Theorem~\ref{thm:nonprojective} are equivalent, with no restriction
on $D\cdot F$.
\end{theorem}
\begin{proof}
Necessity and positive-eigenvalue boundary components are handled
by Lemma~\ref{lem:surface-boundary-tools}. For a connected component
$E$ with zero largest eigenvalue, choose its positive kernel divisor
$A$. If $E$ is vertical, it is a whole fiber, and the punctured-base
construction applies. Otherwise $A\cdot F>0$ excludes vertical
components. The ruling is projective; choose a section $C_0$ with
$C_0^2=-e$. From
$A\equiv aC_0+bF$, $A^2=0$, and $A\cdot E_i=0$ one obtains
\[
 b=ae/2,\qquad E_i\equiv d_i(C_0+(e/2)F),\qquad
 E_i^2=E_i\cdot E_j=0\quad(i\ne j).
\]
Thus $E=C$ is one smooth square-zero curve. For $d=C\cdot F$,
adjunction gives $2g(C)-2=d(2g(B)-2)$, so $C\to B$ is finite
\'etale by Riemann--Hurwitz.

For $d=1$, the complement is an affine-line bundle with degree-zero
linear bundle. The construction of \cite[Proposition~2.1]{KoikeUeda2025}
can be written in unitary flat frames as follows. Choose a smooth
section whose $\bar\partial$ is harmonic. Local antiholomorphic
primitives $g_i$, followed by holomorphic translations of the affine
coordinates, give
\[
 w_i=a_{ij}w_j+c_{ij},\qquad g_i=a_{ij}g_j+c_{ij},
 \qquad |a_{ij}|=1,
\]
with $a_{ij},c_{ij}$ constant. The functions $|w_i-g_i|^2$ agree;
their Levi matrices in $(z,w_i)$ are
$\operatorname{diag}(|(g_i)_{\bar z}|^2,1)$.
They give a smooth psh exhaustion over compact $B$.

For $d\ge2$, a finite \'etale Galois base change splits $C$ into
$d$ disjoint square-zero sections. If $d=2$, the complement is
$L^\times$ with $\deg L=0$, exhausted by
$(\log\|v\|^2)^2$ for a unitary flat metric. If $d\ge3$, three
sections trivialize the ruling. The others are constant maps into
$\PP^1\setminus\{0,1,\infty\}$, so the complement is
$B'\times(\PP^1\setminus Q)$. Pull back an exhaustion of the
punctured fiber. Averaging a nonnegative exhaustion over the finite
deck group preserves properness and descends smoothly. Restrict the
result to a collar at $E$ and glue the boundary collars.
\end{proof}

This criterion uses the actual smooth ruling. Theorem~\ref{thm:divisor-regularization}
gives regularization on arbitrary projective compactifiers, but does
not assert this numerical existence criterion for them.

\subsection{Square-zero curves}\label{subsec:smooth-boundary}
For a smooth square-zero curve, infinite Ueda type is a formal
condition. A unitary flat neighborhood requires convergent defining
functions $w_i=t_{ij}w_j$, $|t_{ij}|=1$.

\begin{proposition}\label{prop:surface-ueda-sufficient}
Let $C\subset S$ be smooth, connected, and square-zero, with $S$
compact K\"ahler. Put $N=N_{C/S}$, let
$\rho:H_1(C,\mathbb Z)\to\Uone$ be its unitary character, and set
$K=\ker[H_1(C,\mathbb Z)\to H_1(S,\mathbb Z)]$.
The complement $S\setminus C$ has a smooth psh exhaustion under
any of the following conditions:
\[
 \begin{gathered}
 C\text{ has finite Ueda type or a unitary flat neighborhood};\\
 N\text{ is torsion or Diophantine};\qquad \rho(K)\text{ is finite};\\
 q(S)\ge g(C);\qquad g(C)=1.
 \end{gathered}
\]
Here Diophantine means
$\operatorname{dist}(N^m,\OO_C)\ge c m^{-\tau}$ for some
$c,\tau>0$ and all $m\ge1$. Only the elliptic-curve condition uses
\cite[Theorem~1.6]{CFSTZ2026}.
\end{proposition}
\begin{proof}
Finite type gives a smooth strictly psh collar
\cite[Theorem~2.5]{CLPT2018}; a flat neighborhood gives
$-\log|w_i|^2$. In infinite type, Ueda's convergence theorem
\cite[Theorem~3]{Ueda1983} handles torsion or Diophantine $N$.
If $\rho(K)$ is finite, some $\rho^k$ factors through the image of
$H_1(C,\mathbb Z)$ in $H_1(S,\mathbb Z)$ and extends to the latter
group because $\Uone$ is divisible. In the remaining non-torsion,
infinite-type case, \cite[Theorem~4.4]{CLPT2018} gives an effective
$E$ disjoint from $C$ and cohomologous to $mC$.
A flat metric on $\OO_S(mC-E)$ and its canonical meromorphic
section $\sigma$ give the collar $-m^{-1}\log\|\sigma\|^2$.

For $q(S)\ge g(C)$, if
$r:H^1(S,\OO_S)\to H^1(C,\OO_C)$ is not injective,
\cite[Proposition~2.7]{CLPT2018} makes a multiple of $C$ a fiber,
and the punctured-base construction applies. Otherwise $r$ is an
isomorphism. Hodge theory gives an isomorphism on rational first
homology; since $H_1(C,\mathbb Z)$ is torsion-free, $K=0$.
Finally, if $g(C)=1$, \cite[Theorem~1.6]{CFSTZ2026} rules out
infinite type without a unitary flat neighborhood. Apply
Lemma~\ref{lem:surface-boundary-tools}\textup{(ii)} in all cases.
\end{proof}

\section{Holomorphic convexity for curve fiberings}\label{sec:curve-fibers}
Complex spaces in this section are reduced, Hausdorff, and countable at
infinity; smooth functions are understood in local embeddings. Fiber
components refer to reduced supports.

\begin{theorem}\label{thm:curve-fiber-convexity}
Let $X$ have pure dimension $n\ge2$ and admit a continuous real-valued
psh exhaustion $\rho$. If a holomorphic map
\[
 F:X\longrightarrow\C^{n-1}
 \quad\text{satisfies}\quad
 \dim F^{-1}(z)\le1\qquad(z\in F(X)),
\]
then $X$ is holomorphically convex and admits a smooth psh exhaustion.
\end{theorem}

For smooth $X$, the convexity argument follows
\cite[Corollary~4.2 and pp.~463--464]{VijituForum2022}.
The cycle argument below allows arbitrary singularities of $X$.

\subsection{Compact fiber components}
The analytic dimension theorem gives
\[
 \dim_x F^{-1}(F(x))\ge n-(n-1)=1.
\]
Applied on each local irreducible branch, it shows that every nonempty
fiber is pure one-dimensional. Since the base is normal, $F$ has
analytic weighted fiber cycles $\mathcal F_z$ satisfying
\[
 |\mathcal F_z|=F^{-1}(z),
 \qquad\mathcal F_z=\sum_jm_j(z)[C_j(z)],\qquad m_j(z)\in\mathbb N_{>0};
\]
see \cite[Chapter~I, \S1(b), p.~816]{BarletMagnusson1998}.
This assertion concerns geometric cycles and requires neither
properness nor flatness of the local rings.

\begin{lemma}[Uniform area]\label{curve:lem:area}
Fix a Hermitian metric $h$ on $X$. Given compact $L\subset X$ and
$z_0\in\C^{n-1}$, there are a neighborhood $B$ of $z_0$ and $A>0$
such that
\[
 \int_C h\le A
\]
for every compact irreducible component $C\subset L$ of a fiber over $B$.
\end{lemma}
\begin{proof}
Choose $\chi\in C_c^\infty(X)$ with $\chi\ge0$ and $\chi=1$ near $L$.
Integration of a compactly supported continuous form is continuous
on an analytic family of cycles
\cite[Chapter~IV, Proposition~2.3.1]{BarletMagnusson2014}.
Consequently, after shrinking $B$,
\[
 \int_C h=\int_{[C]}\chi h
 \le\int_{\mathcal F_z}\chi h\le A\qquad(z\in B).
\]
\end{proof}

\begin{lemma}[Stability]\label{curve:lem:stability}
Let $K\subset X$ be compact. If every irreducible component of
$F^{-1}(z_0)$ meeting $K$ is noncompact, the same holds for all
fibers over a neighborhood of $z_0$.
\end{lemma}
\begin{proof}
Otherwise choose $z_j\to z_0$, compact irreducible components
$C_j\subset F^{-1}(z_j)$, and $x_j\in C_j\cap K$, with $x_j\to x\in K$.
The maximum principle on the normalization of $C_j$ makes $\rho|_{C_j}$
constant. Hence
\[
 C_j\subset L:=\{\rho\le\max_K\rho\}\Subset X.
\]
Lemma~\ref{curve:lem:area} bounds their areas. Bishop compactness for
cycles on complex spaces \cite[Theorem~1]{Barlet1999} gives, after
passing to a subsequence, a compact effective one-cycle
\[
 [C_j]\longrightarrow T,\qquad x\in|T|\subset F^{-1}(z_0).
\]
Indeed, upper continuity of compact supports prevents the limit from
losing $x$, and local convergence of integration currents puts its
support in the central fiber. An irreducible component of $|T|$
through $x$ is a compact analytic curve in a pure one-dimensional
fiber, hence a full irreducible component of that fiber. This
contradicts the hypothesis. The cases $K=\varnothing$ or
$F^{-1}(z_0)\cap K=\varnothing$ are immediate.
\end{proof}

\subsection{Convexity and reduction}
We use three complex-space criteria from
\cite[Propositions~2.2, 2.4 and Corollary~2.5]{VijituForum2022}:
\begin{enumerate}[label=\textup{(\roman*)},leftmargin=2.5em]
\item A complex space with a continuous psh function $u$ is
holomorphically convex if and only if every $\{u<t\}$ is.
\item A pseudoconvex space locally holomorphically convex over a
Stein base is holomorphically convex.
\item If $\pi:M\to S$ is holomorphic, $S$ is Stein, and
$D\subset M$ is pseudoconvex, then $D$ is holomorphically convex
provided $\overline D^{\,M}\to S$ is proper and a neighborhood
of $\partial_M D$ has Stein reduced fibers.
\end{enumerate}
In \textup{(iii)}, the auxiliary smooth function $h$ constructed in the
cited proof must be supplemented by an exhaustion $\eta_D$ of $D$.
Over a sufficiently small relatively compact base neighborhood,
$h$ is bounded on $\overline D$ and strictly psh in a boundary collar.
Thus $\eta_D+h$ is a vertical exhaustion and is strictly psh above
a constant: outside the collar it is bounded on a compact subset of
$D$. This supplies the required $1$-convex map, to which
\cite[Theorem~2.3]{VijituForum2022} applies.

\begin{lemma}\label{curve:lem:neighborhood}
For every $\lambda\in\R$ and $z_0\in\C^{n-1}$, there are $c>\lambda$
and a ball $B=B_r(z_0)$ such that
\[
 D=F^{-1}(B)\cap\{\rho<c\}
 \quad\text{and}\quad
 F^{-1}(B)\cap\{\rho<\lambda\}
\]
are holomorphically convex.
\end{lemma}
\begin{proof}
Let $S_0$ be the set of constant $\rho$-values on compact irreducible
components of $F^{-1}(z_0)$. Local finiteness of fiber components
and compactness of $\{\rho\le R\}$ imply
\[
 \#\bigl(S_0\cap(-\infty,R]\bigr)<\infty\qquad(R\in\R).
\]
Choose $\lambda<a<c<b$ with $[a,b]\cap S_0=\varnothing$.
Lemma~\ref{curve:lem:stability}, applied to
$K=\{a\le\rho\le b\}$, provides a ball $B=B_r(z_0)$ such that no
compact irreducible component of a fiber over $B$ meets $K$.
Set
\[
 M=F^{-1}(B),\qquad D=M\cap\{\rho<c\},\qquad
 U=M\cap\{a<\rho<b\}.
\]
Then $\partial_M D\subset U$. A compact irreducible component of a
reduced $U$-fiber would be a compact analytic curve in $X$, hence
a full component of the corresponding full fiber meeting $K$.
Thus all reduced $U$-fibers are Stein, by the Stein criterion for
curves without compact irreducible components
\cite[\S3, p.~1663]{Vajaitu2022}.

The function
\[
 \eta=\max\left\{\frac1{c-\rho},
              \frac1{r^2-\|F-z_0\|^2}\right\}
\]
is a continuous psh exhaustion of $D$: for $t>0$,
\[
 \{\eta\le t\}
 =\{\rho\le c-1/t\}\cap
 F^{-1}\{\|z-z_0\|^2\le r^2-1/t\}\Subset D.
\]
Moreover, $\overline D^{\,M}\to B$ is proper, since its inverse
image over compact $T\subset B$ is closed in
$\{\rho\le c\}\cap F^{-1}(T)$.
Criterion~\textup{(iii)} makes $D$ holomorphically convex;
\textup{(i)} applied to $\rho|_D$ gives the second assertion.
\end{proof}

\begin{proof}[Proof of Theorem~\ref{thm:curve-fiber-convexity}]
For $X(\lambda)=\{\rho<\lambda\}$, Lemma~\ref{curve:lem:neighborhood}
makes $F|_{X(\lambda)}$ locally holomorphically convex over
$\C^{n-1}$. The function $(\lambda-\rho)^{-1}$ exhausts $X(\lambda)$,
so criterion~\textup{(ii)} makes $X(\lambda)$ holomorphically convex.
Criterion~\textup{(i)} now gives holomorphic convexity of $X$.
For its proper Remmert reduction $q:X\to Z$, pull back a smooth
strictly psh exhaustion $u$ of the Stein space $Z$:
\[
 u\circ q\in C^\infty(X)\cap\operatorname{PSH}(X),\qquad
 \{u\circ q\le t\}=q^{-1}(\{u\le t\})\Subset X.
\]
This also proves Theorem~\ref{thm:curve-fiber-introduction}.
\end{proof}

\begin{corollary}[Simultaneous contraction]\label{cor:curve-contraction}
The Remmert reduction $q:X\to Z$ has connected fibers,
$q_*\OO_X=\OO_Z$, and $F=\widetilde F\circ q$ for a holomorphic
$\widetilde F:Z\to\C^{n-1}$. Its nontrivial fibers are exactly the
maximal connected unions of compact irreducible fiber curves of $F$.
\end{corollary}
\begin{proof}
The sheaf identity makes every component of $F$ descend.
A nontrivial $q$-fiber is therefore a compact connected analytic
subset of a one-dimensional $F$-fiber, hence a finite connected
union of its compact irreducible components. Conversely, every
compact irreducible curve maps to a point of $Z$. A connected union
of such curves has constant $\rho$, lies in a compact sublevel,
and is finite by local finiteness of the components of its $F$-fiber.
It is consequently contained in a $q$-fiber, proving maximality.
\end{proof}

\begin{thebibliography}{99}
\raggedright

\bibitem{Koike2021}
T.~Koike,
\emph{Linearization of transition functions of a semi-positive line
bundle along a certain submanifold},
Ann. Inst. Fourier \textbf{71} (2021), no.~5, 2237--2271.
\href{https://doi.org/10.5802/aif.3439}{doi:10.5802/aif.3439}.

\bibitem{Wu2026}
X.~Wu,
\emph{Generalized Ueda Obstruction Classes and Applications to
Non-semi-positivity of Line Bundles},
preprint (2026), arXiv:2606.22821v1.
\href{https://arxiv.org/abs/2606.22821v1}{arXiv:2606.22821}.

\bibitem{Zhang2025}
Y.~Zhang,
\emph{On the Non-Semipositivity of a Nef and Big Line Bundle on
Grauert's Example},
preprint (2025), arXiv:2512.23446v1.
\href{https://arxiv.org/abs/2512.23446v1}{arXiv:2512.23446}.

\bibitem{Vajaitu2022}
V.~V\^{\i}j\^{\i}itu,
\emph{Holomorphic convexity and pseudoconvex surfaces},
Ann. Sc. Norm. Super. Pisa Cl. Sci. (5) \textbf{23} (2022),
no.~4, 1659--1678.
\href{https://doi.org/10.2422/2036-2145.202011_018}{doi:10.2422/2036-2145.202011\_018}.
\url{https://journals.sns.it/index.php/annaliscienze/article/view/5700}.

\bibitem{MongodiSlodkowski2020}
S.~Mongodi and Z.~Slodkowski,
\emph{Domains with a continuous exhaustion in weakly complete surfaces},
Math. Z. \textbf{296} (2020), nos.~3--4, 1011--1019.
\href{https://doi.org/10.1007/s00209-020-02466-z}{doi:10.1007/s00209-020-02466-z}.
\url{https://arxiv.org/abs/1904.03497}.

\bibitem{Grauert1962}
H.~Grauert,
\emph{\"Uber Modifikationen und exzeptionelle analytische Mengen},
Math. Ann. \textbf{146} (1962), 331--368.
\href{https://doi.org/10.1007/BF01441136}{doi:10.1007/BF01441136}.

\bibitem{FN2005}
Y.~Fujimoto and N.~Nakayama,
\emph{Compact complex surfaces admitting non-trivial surjective endomorphisms},
Tohoku Math. J. (2) \textbf{57} (2005), no.~3, 395--426.
\href{https://doi.org/10.2748/tmj/1128703004}{doi:10.2748/tmj/1128703004}.
\href{https://www.kurims.kyoto-u.ac.jp/preprint/file/RIMS1431.pdf}{Author preprint, Section~2 and Lemma~2.1}.

\bibitem{Enoki}
I.~Enoki, \emph{Surfaces of class $\mathrm{VII}_0$ with curves}, Tohoku Math. J. (2)
\textbf{33} (1981), no.~4, 453--492.
\href{https://doi.org/10.2748/tmj/1178229349}{doi:10.2748/tmj/1178229349}.
\href{https://www.jstage.jst.go.jp/article/tmj1949/33/4/33_4_453/_pdf/-char/en}{Original article, pp.453--454}.

\bibitem{Ueda1983}
T.~Ueda,
\emph{On the neighborhood of a compact complex curve with topologically
trivial normal bundle},
J. Math. Kyoto Univ. \textbf{22} (1982/83), no.~4, 583--607.
\href{https://doi.org/10.1215/kjm/1250521670}{doi:10.1215/kjm/1250521670}.

\bibitem{CLPT2018}
B.~Claudon, F.~Loray, J.~V.~Pereira and F.~Touzet,
\emph{Compact leaves of codimension one holomorphic foliations on
projective manifolds},
Ann. Sci. \'Ecole Norm. Sup. (4) \textbf{51} (2018), no.~6, 1457--1506.
\href{https://doi.org/10.24033/asens.2379}{doi:10.24033/asens.2379}.
\href{https://arxiv.org/abs/1512.06623}{Author preprint, Theorems~1.3, 2.5, 4.4 and Proposition~2.7}.

\bibitem{CFSTZ2026}
Y.~Chen, S.~Filip, S.~Sun, V.~Tosatti and J.~Zhang,
\emph{Compact K\"ahler surfaces with semipositive anticanonical bundle},
preprint (2026), arXiv:2609.26716v1, submitted September~22, 2026.
\href{https://arxiv.org/abs/2609.26716v1}{arXiv:2609.26716}.

\bibitem{KoikeUeda2025}
T.~Koike and T.~Ueda,
\emph{Plurisubharmonic functions on affine line bundles over compact
K\"ahler manifolds},
Publ. Res. Inst. Math. Sci. \textbf{61} (2025), no.~1, 139--152.
\href{https://doi.org/10.4171/PRIMS/61-1-3}{doi:10.4171/PRIMS/61-1-3}.

\bibitem{Barlet2015}
D.~Barlet,
\emph{Some examples due to H.~Hironaka},
preprint (2015), arXiv:1501.01129v2, Section~6.
\href{https://arxiv.org/abs/1501.01129}{arXiv:1501.01129}.

\bibitem{StacksHironaka}
The Stacks Project Authors,
\emph{The Stacks Project},
Lemma~110.67.1, Tag~08J1, ``Being projective is not local on the base.''
\url{https://stacks.math.columbia.edu/tag/08J1}.

\bibitem{DNS22}
T.-C.~Dinh, V.-A.~Nguy\^en, and N.~Sibony,
\emph{Unique ergodicity for foliations on compact K\"ahler surfaces},
Duke Math. J. \textbf{171} (2022), no.~13, 2627--2698.
\href{https://doi.org/10.1215/00127094-2022-0044}{doi:10.1215/00127094-2022-0044}.
Author version: \href{https://arxiv.org/abs/1811.07450v2}{arXiv:1811.07450v2}.
The numbered equations and Appendix~B cited here may be found in this
author version.

\bibitem{Wlodarczyk05}
J.~W{\l}odarczyk,
\emph{Simple Hironaka resolution in characteristic zero},
J. Amer. Math. Soc. \textbf{18} (2005), no.~4, 779--822.
\href{https://doi.org/10.1090/S0894-0347-05-00493-5}{doi:10.1090/S0894-0347-05-00493-5}.
Author version: \href{https://arxiv.org/abs/math/0401401}{arXiv:math/0401401}.

\bibitem{VijituForum2022}
V.~V\^{\i}j\^{\i}itu,
\emph{Holomorphic convexity of pseudoconvex spaces in terms of
the rank of structural sheaf},
Forum Math. \textbf{34} (2022), no.~2, 455--468.
\href{https://doi.org/10.1515/forum-2021-0262}{doi:10.1515/forum-2021-0262}.

\bibitem{BarletMagnusson1998}
D.~Barlet and J.~Magn\'usson,
\emph{Int\'egration de classes de cohomologie m\'eromorphes et
diviseurs d'incidence},
Ann. Sci. \'Ecole Norm. Sup. (4) \textbf{31} (1998),
no.~6, 811--842.
\href{https://doi.org/10.1016/S0012-9593(99)80004-4}
{doi:10.1016/S0012-9593(99)80004-4}.
\url{https://www.numdam.org/item/ASENS_1998_4_31_6_811_0.pdf}.

\bibitem{BarletMagnusson2014}
D.~Barlet and J.~Magn\'usson,
\emph{Cycles analytiques complexes I: th\'eor\`emes de
pr\'eparation des cycles},
Cours Sp\'ecialis\'es, vol.~22,
Soci\'et\'e Math\'ematique de France, Paris, 2014.

\bibitem{Barlet1999}
D.~Barlet, \emph{How to use the cycle space in complex geometry},
in \emph{Several Complex Variables},
M.~Schneider and Y.-T.~Siu (eds.),
MSRI Publications, vol.~37, Cambridge University Press, 1999,
pp.~25--42.
\url{https://library.slmath.org/books/Book37/files/barlet.pdf}.

\end{thebibliography}
\end{document}
